% \setlength{\columnsep}{12pt}
% espaço entre colunas %\setlength{\columnseprule}{0.2pt} % linha vertical entre colunas

%-------------------------------------------------------------------------
% Modelo Trabalho PPGE-UFPB/Estatística
% Normas da ABNT
% Provedor: Laboratório de Economia e Modelagem Aplicada -- LEMA
%-------------------------------------------------------------------------

\documentclass[12pt, openright,oneside, a4paper, english, brazil]{report}

% ----------------------------
% Pacotes básicos 
% ----------------------------
% Fonte e codificação (acentuação)
\usepackage{lmodern}

\renewcommand{\sfdefault}{ppl}
\usepackage[T1]{fontenc}	
\usepackage[utf8]{inputenc}	
\usepackage[brazil]{babel}
\usepackage{comment} % para esconder a tabela Org do compilador
% Referências
%\usepackage[brazilian,hyperpageref]{}	
%\usepackage[alf]{abntex2cite}			
\usepackage[colon]{natbib}
\bibliographystyle{apalike}

%Texto aleatório
\usepackage{lipsum} 

% Tabelas e Figuras
\usepackage{ctable}             
\usepackage{multirow}
\usepackage{array}
\usepackage{float}              
\usepackage{longtable}          
\usepackage{xcolor,colortbl}    
\usepackage{booktabs}           
\usepackage{graphicx}			
\usepackage{subfig}             
\usepackage{caption}            
\usepackage{epstopdf}           

% Equações e símbolos
\usepackage{amsmath}            
\usepackage{nomencl} 
\usepackage{hyperref, cleveref}  

% Texto
\usepackage{indentfirst}		
\usepackage{color}				
\usepackage{microtype} 			
\usepackage{lastpage}			
% ----------------------------------------------------------
% Configurações do texto
% ----------------------------------------------------------
\usepackage{multicol}
\usepackage{array}
\newcolumntype{L}[1]{>{\raggedright\let\newline\\\arraybackslash\hspace{0pt}}m{#1}}
\newcolumntype{C}[1]{>{\centering\let\newline\\\arraybackslash\hspace{0pt}}m{#1}}
\newcolumntype{R}[1]{>{\raggedleft\let\newline\\\arraybackslash\hspace{0pt}}m{#1}}

\makeatletter
%\hypersetup{pdftitle={\@title}, pdfauthor={\@author}}
\makeatother

\def\preambulo{Trabalho elaborado em cumprimento às exigências da disciplina \disciplina, ministrada pela professor Dr. \professor.}
%--------------------------------------------------------------

\setlength{\parindent}{1.25cm}
\setlength{\parskip}{0.1cm} 

\usepackage{titlesec}
\usepackage{lipsum}

\makeindex




\title{Exercício Estudo de Caso}
\def\autor{Luiz Fernando Ferraz\\Oscar Sérgio Monteiro}
\def\disciplina{SC334 - Avaliação de Desempenho Organizacional, turma b}
\def\professor{Antônio César Pitela}
\begin{document}
\makeatletter
\begin{titlepage}
    \begin{figure}[H]
    \vspace{-3cm}
    \centering
    \includegraphics[height=3cm,keepaspectratio]{config/coverpage}\hfill
    \par\vspace{1cm}
    \end{figure}
    \begin{figure}[H]
    \vspace{-2cm}
    \centering
    \includegraphics[height=3cm,keepaspectratio]{config/contabeis.png}
    \par\vspace{1cm}
    \end{figure}
     
    \vspace{2cm}
    
	\begin{center}
	    {\huge \bfseries \@title}
	    \begin{flushright}
		\begin{minipage}{8cm}
			\preambulo
		\end{minipage}
	\end{flushright}
	\vspace{3cm}
	    \normalsize
	    \autor
	\end{center}
	
\vfill
\centering
{\large Curitiba-PR, \the\year}

\end{titlepage}
\makeatother


\pdfbookmark[0]{\contentsname}{toc}
\tableofcontents
\newpage


\titleformat{\chapter}[display]{\normalfont\bfseries}{}{0pt}{\Huge  \thechapter. }
\newcommand{\doubletoprule}{%
  \specialrule{.1em}{2pt}{2pt}%
  \specialrule{.1em}{0pt}{0pt}%
}

\chapter{Caso Conecta Serviços Empresariais}
\section{Atividades}
\subsection{Questão 6.1 - Indicadores}
Calcule, para 2024 e 2025, o crescimento da receita, as margens bruta, operacional e líquida, o ROA e o ROE. Compare 2025 com as metas.
\begin{tabular}{p{5cm} l l l l }                                                     
\toprule                                                                       
\textbf{Indicador} & \textbf{2024} &	\textbf{2025} &	 \textbf{Meta 2025} & \textbf{Conclusão} \\
\midrule
Crescimento da receita & -  & 18,00\%	& $\geq$  15,00\% & Ok \\
\midrule
Margem bruta &	37,50\% & 31,85\% & $\geq$ 36,00\% & $\downarrow$   \\
\midrule
Margem operacional & 15,00\% & 	10,31\% & $\geq$ 14,00\% & $\downarrow$  \\
\midrule
Margem líquida	& 9,58\% & 6,21\% &  $\geq$  9,00\% &	 $\downarrow$  \\
\midrule
ROA & 16,43\% &	10,48\% & $\geq$ 15,00\% &	 $\downarrow$  \\
\midrule
ROE & 28,75\% &	20,23\%	& $\geq$ 25,00\% & $\downarrow$ \\
 \bottomrule                                                                    
 \end{tabular}\\                                                                  
A exceção do crescimento das receitas todas os demais indicadores ficaram abaixo da meta\\ 
\subsection{Questão 6.2 - Comparação das carteiras}
Calcule, para as carteiras Essencial e Premium, a margem de contribuição e a receita média por cliente. O maior preço médio da Premium está se convertendo em melhor desempenho?\\
 \begin{tabular}{p{2cm} l l l}                                                     
 \toprule                                                                       
 \textbf{Carteira} & \textbf{Descrição} & \textbf{Referência} & \textbf{Indicador} \\
 \midrule
 Essencial & Margem de Contribuição & (8.100 – 5.270) / 8.100 * 100  & 34,94\% \\
 \midrule                                                                       
 Essencial & Receita média por cliente  & 8.100 / 130  & R\$ 62.307,69  \\                                        
 \midrule                                                                       
 Premium &  Margem de Contribuição & (6.060 – 4.300) / 6.060 * 100  & 29,04\% \\
 \midrule                                                                       
 Premium &  Receita média por cliente  & 6.060 / 48 & R\$ 126.250,00  \\ 
 \bottomrule                                                                    
 \end{tabular}
 Observa-se que a maior receita por cliente não significa necessariamente uma melhor margem de contribuição\\
\subsection{Questão 6.3}
Classifique os indicadores do caso nas perspectivas financeira, clientes, processos internos e aprendizado e crescimento.
\subsubsection{Indicadores não financeiros}
\hspace*{-2cm}                                                                 
 \begin{tabular}{p{4cm} l }                                                     
 \toprule                                                                       
   \textbf{Indicador} & \textbf{Categoria} \\
 \midrule
Crescimento da receita & financeiro \\
\midrule
Margem bruta & financeiro \\
\midrule
Margem operacional & financeiro \\
\midrule
Margem líquida & financeiro \\
\midrule
ROA & financeiro \\
\midrule
ROE & financeiro \\
\midrule
Margem de contribuição & financeiro \\
\midrule
Receita média por cliente & financeiro \\
\midrule
Clientes ativos & clientes \\
\midrule
Renovação de contratos & clientes \\
\midrule
Churn de clientes & clientes \\
\midrule
NPS & clientes \\
\midrule
Atendimentos no SLA & processo \\
\midrule
Resolução no primeiro contato & processo \\
\midrule
Tempo médio de resolução & processo \\
\midrule
Chamados reabertos & processo \\
\midrule
Utilização faturável & processo \\
\midrule
Horas extras mensais por empregado & aprendizado \\
\midrule
Rotatividade de pessoal  & aprendizado \\
\midrule
Absenteísmo  & aprendizado \\
\midrule
Treinamento por empregado  & aprendizado \\
\midrule
Engajamento & aprendizado \\
 \bottomrule                                                                    
 \end{tabular}\\
 Face a tantos indicadores fora da meta uma analise mais criteriosa deve ser realizada, entendemos que não se trata de um problema especifico de cada setor e sim problemas de ordem sistêmica. De qualquer forma recomendamos iniciar pelo indicador de Horas de Retrabalho; para um acréscimo de 20\% no número de peças produzidas o tempo de retrabalho foi acrecido em 83\%.
\subsection{Questão 6.4}
Identifique pelo menos quatro relações de causa e efeito. Diferencie as relações apoiadas pelos dados das hipóteses que ainda exigem investigação.
\textbf{Resposta}\\
\begin{itemize}
\item Treinamento e Tempo médio de resolução\\
  A meta de 24 horas de treinamento por funcionário indica a necessidade de atualização constante, esta atualização esta diretamente relacionada ao desempenho dos colaboradores no atendimento ao clientes. A diminuição para o patamar de 10 horas em 2025 podem justificar o aumento do indicador de ``Tempo médio de resolução'' que saiu de 7,5 para 11,2. Estes valores indicam que as horas disponíveis para treinamento são compensadas pelo tempo gasto na solução de problemas.\\
\item Treinamento e Atendimento de SLA\\
  Assim como o item anteriores o não atingimento das metas de treinamento são muito custosas para as empresas prestadoras de serviço. O atingimento no nível de atendimento é diretamente relacionada com a eficácia e eficiência da prestação do serviço.\\
\item Outros custos de serviços vs Rotatividade de Pessoal\\
  A rotatividade do pessoal tem uma relação direta com o aumento de Outros Custos de serviços , em especial os custos de contratação e de treinamento.\\
\item Custos diretos de Pessoal vs Horas extras mensais por empregado\\
Não há dúvida que a relação e direta, cabe observar as causas do aumento do numero de horas extras por funcionário. \\
\end{itemize}
\subsection{Questão 6.5}
  A utilização faturável de 88\% deve ser interpretada como eficiência superior? Discuta o indicador em conjunto com horas extras, rotatividade, SLA, reabertura e tempo de resolução.\\
 \textbf{Resposta}\\
  Talvez mais apropriado seria dizer que é um indicador de eficácia , a medida que mais disponibilidade de tempo para faturamento indica que os funcionários estão mais tempo disponíveis para fazer o que traz receia para a empresa e menos tempo executando tarefas secundárias. Assim por isso a meta para o indicador é de 95\%. No caso em questão observa-se inconsistência desse indicador com outros também relacionados a eficácia e eficiência. O aumento no numero de horas extras aumentado significa que parte das horas disponíveis para faturamento não aconteceu no horário normal de trabalho, em tese horas extras só deveriam ocorrer com níveis mais altos de utilização faturável. Outro indicativo de inconsistência é a diminuição das horas de treinamento, poderia indicar que horas deveriam estar em treinamento foram disponibilizadas para faturamento, mas outra vez essa situação só deveria ocorrer se a Utilização faturável estivesse acima dos 95\%.
\subsection{Questão 6.6}
Compare as carteiras. A empresa deve acelerar, manter ou suspender temporariamente as vendas da Premium? Defenda uma alternativa e apresente os riscos da decisão.\\
\textbf{Resposta}\\
A sustentabilidade de uma empresa também depende da capacidade de inovação, considerando que o atendimento premium é uma inovação de processo entendo que é prematuro descontinuar a nova atividade, mesmo com as evidências de piora nos níveis gerais de atendimento e na rentabilidade da empresa. No memento a prudencia recomendaria desacelerar a implantação de inovação, identificar as causas dos problemas, e tentar sanar os problemas. Não existem crescimento, inovação ou mudança sem esforço e risco. A revisão do plano de implantação da nova modalidade de atendimento deve ser revisado, em especial a gestão de riscos.\\
\subsection{Questão 6.7}
Construa um Balanced Scorecard enxuto, com dois indicadores por perspectiva. Para cada indicador,informe meta, periodicidade e responsável.
\hspace*{-2cm}
\begin{table}[htbp]                                                            
 \centering                                                                     
  \begin{tabular}{|p{3cm}|p{3cm}|p{2cm}|p{2cm}|p{3cm}|}                      
 \hline                                                                         
 \textbf{Perspectiva} & \textbf{Indicador} & \textbf{Meta} &                    
 \textbf{Periodic.} & \textbf{Resp.} \\                               
 \hline                                                                         
 Financeira & Receita por cliente ativo & $\geq$ R\$ 80 mil/ano & Mensal &      
 Financeiro / Comercial \\                                                        
 \hline                                                                         
 Financeira & Índice de custos diretos sobre a receita & $\leq$ 49\% & Mensal & 
 Financeiro \\                                                                  
 \hline                                                                         
 Clientes & Taxa de retenção da carteira Premium & $\geq$ 92\% & Mensal &       
 Comercial / Relacionamento \\                                                    
 \hline                                                                         
 Clientes & Índice de reclamações por 100 clientes & $\leq$ 3 & Mensal &        
 Atendimento / Qualidade \\                                                       
 \hline                                                                         
 Processos internos & Tempo de resposta inicial ao cliente & $\leq$ 2 horas &   
 Semanal & Operações \\                                                         
 \hline                                                                         
 Processos internos & Percentual de chamados concluídos sem escalonamento &     
 $\geq$ 85\% & Semanal & Operações/Suporte \\                                   
 \hline                                                                         
 Aprendizado e crescimento & Índice de cobertura de capacitação crítica &       
 $\geq$ 90\% & Mensal & RH / Treinamento \\                                       
 \hline                                                                         
 Aprendizado e crescimento & Clima organizacional da equipe operacional &       
 $\geq$ 75 & Trimestral & RH \\                                                 
 \hline                                                                         
 \end{tabular}                                                                  
 \end{table}                                                                    
\subsection{Questão 6.8}
Proponha um plano com quatro ações prioritárias e elabore um parecer gerencial de 12 a 18 linhas.\\
\textbf{Resposta}\\
\begin{itemize}
\item Revisar o plano de implantação da nova carteira , comparar os resultados esperados com o resultado obtido.\\
\item Fazer um nova plano de gerenciamento de riscos.\\
\item Segregar os indicadores por carteira.\\
\item Especializar as equipes por carteira.\\
\end{itemize}\\
Observa-se uma deterioração nos indicadores de desempenho de 2024 para 2025, a principal causa esta na implantação de uma nova carteira de serviço. Esta deterioração parece indicar que o esforço para a implantação do novo serviço foi subestimado. Justamente porque as inovações são muito importantes para a sustentabilidade das empresas, devem ser muito bem planejadas. Não se pode conceber a implantação de uma nova carteira de serviços sem um exaustivo plano de ação e especialmente um plano piloto que permita mitigar as incertezas da inovação. A situação posta requer extremo cuidado, visto que já foi criada expectativa nos clientes de uma nova forma de atendimento, desta forma não se pode pedir aos clientes que já optaram pelo novo serviço que esperem até a empresa se organizar para o atendimento e resolver os problemas que surgiram. A solução dos problemas deverá ocorrer ``ao vivo'' ou seja sem possibilidade de ensaio. Isso coloca a empresa em uma situação de alto risco, desta forma todas as ações devem ser cautelosas, e uma atenção especial deve ser dada a operação tradicional visto que é ele que vai manter a empresa nesta fase de implantação do novo serviço. É importante isolar os problemas da nova operação para não contaminar a operação onde a empresa tem experiência.\\
\end{document}
